\documentclass[a4paper,11pt]{article}
\usepackage{amsmath,amsthm,amsfonts,amssymb,amscd,amstext,vmargin,graphics,graphicx,tabularx,multicol} 
\usepackage[francais]{babel}
\usepackage[utf8]{inputenc}  
\usepackage[T1]{fontenc} 
\usepackage{pstricks-add,tikz,tkz-tab,variations}
\usepackage[autolanguage,np]{numprint} 

\setmarginsrb{1.5cm}{0.5cm}{1cm}{0.5cm}{0cm}{0cm}{0cm}{0cm} %Gauche, haut, droite, haut
\newcounter{numexo}
\newcommand{\exo}[1]{\stepcounter{numexo}\noindent{\bf Exercice~\thenumexo} : \marginpar{\hfill /#1}}
\reversemarginpar


\newcounter{enumtabi}
\newcounter{enumtaba}
\newcommand{\q}{\stepcounter{enumtabi} \theenumtabi.  }
\newcommand{\qa}{\stepcounter{enumtaba} (\alph{enumtaba}) }
\newcommand{\initq}{\setcounter{enumtabi}{0}}
\newcommand{\initqa}{\setcounter{enumtaba}{0}}

\newcommand{\be}{\begin{enumerate}}
\newcommand{\ee}{\end{enumerate}}
\newcommand{\bi}{\begin{itemize}}
\newcommand{\ei}{\end{itemize}}
\newcommand{\bp}{\begin{pspicture*}}
\newcommand{\ep}{\end{pspicture*}}
\newcommand{\bt}{\begin{tabular}}
\newcommand{\et}{\end{tabular}}
\renewcommand{\tabularxcolumn}[1]{>{\centering}m{#1}} %(colonne m{} centrée, au lieu de p par défault) 
\newcommand{\tnl}{\tabularnewline}

\newcommand{\trait}{\noindent \rule{\linewidth}{0.2mm}}
\newcommand{\hs}[1]{\hspace{#1}}
\newcommand{\vs}[1]{\vspace{#1}}

\newcommand{\N}{\mathbb{N}}
\newcommand{\Z}{\mathbb{Z}}
\newcommand{\R}{\mathbb{R}}
\newcommand{\C}{\mathbb{C}}
\newcommand{\Dcal}{\mathcal{D}}
\newcommand{\Ccal}{\mathcal{C}}
\newcommand{\mc}{\mathcal}

\newcommand{\vect}[1]{\overrightarrow{#1}}
\newcommand{\ds}{\displaystyle}
\newcommand{\eq}{\quad \Leftrightarrow \quad}
\newcommand{\vecti}{\vec{\imath}}
\newcommand{\vectj}{\vec{\jmath}}
\newcommand{\Oij}{(O;\vec{\imath}, \vec{\jmath})}
\newcommand{\OIJ}{(O;I,J)}


\newcommand{\reponse}[1][1]{%
\multido{}{#1}{\makebox[\linewidth]{\rule[0pt]{0pt}{20pt}\dotfill}
}}

\newcommand{\titre}[5] 
% #1: titre #2: haut gauche #3: bas gauche #4: haut droite #5: bas droite
{
\noindent #2 \hfill #4 \\
#3 \hfill #5

\vspace{-1.6cm}

\begin{center}\rule{6cm}{0.5mm}\end{center}
\vspace{0.2cm}
\begin{center}{\large{\textbf{#1}}}\end{center}
\begin{center}\rule{6cm}{0.5mm}\end{center}
}



\begin{document}
\pagestyle{empty}
\titre{Le petit truc en plus}{}{}{4ème}{}

\vspace*{1cm}

{\Large \textbf{A rendre avant le 13 décembre !!}}

\vspace*{1cm}
La construction d'un sangaku est à réaliser sur une feuille blanche.\\

\includegraphics[scale=1]{sangaku.eps} \\

Un programme de construction :\\
 
\bi \item Placer un point O.\\
\item Construire le cercle de centre O et de diamètre 12 cm.\\
\item Placer un point A sur le cercle et tracer le segment [OA].\\
\item Construire le point B image du point A par la rotation de centre O et d'angle 60\degre (sens anti-horaire) et tracer le segment [OB].\\
\item Construire le point C image du point B par la rotation de centre O et d’angle 60\degre (sens anti-horaire)et tracer le segment [OC].\\
\item Même consigne pour les points D, E, F.\\
\item Construire le cercle de centre O et de rayon 2 cm, il coupe le segment [OA] en A’, le segment [OB] en B’, … , le segment [OF] en F’.\\
\item Construire les cercles de diamètres [AA’], [BB’], [CC’], [DD’], [EE’], [FF’].\\
\ei

\vspace*{1cm}


Temps passé à faire réaliser la figure : . . . .\\

\vspace*{1cm}

\begin{flushright}
\begin{tabular}{|p{9cm}|}
\hline 
\textbf{Critères d'évaluation :} \\ 
\hline 
Je sais faire une figure complète\\
Je sais construire une figure avec précision\\
Je sais construire une figure de façon soignée\\
Je comprends les consignes et le vocabulaire mathématique \\
\hline 
\end{tabular} 

\end{flushright}

\end{document}
